\chapter{Podsumowanie}

%% wyszczególnić które cele projektu zostały wykonane

Okres projektowania i~konstruowania wykonanego w~ramach projektu systemu był dla
autora niewątplie czasem niezwykle kształcącym i~rozwijającym. System,
konstruowany i~projektowany od fundamentów, pozwolił autorowi na dokładną
obserwację wpływu poszczególnej decyzji projektowej na finalny sposób działania
układu. Aby system mógł wykonać swoje zadanie w najprostszej formie, każdy
element musiał idealnie ze sobą współgrać, więc nawet najmniejsza, błędna
decyzja projektowa wykonana w~fazach początkowych, mogła finalnie spowodować
wiele problemów w fazach finalnych, powodując nawet zakończenie projektu
fiaskiem. Poprawne działanie systemu dla sygnałów wąskopasmowych, które
zaobserowowano badając działanie dla poszczególnych sygnałów symulujących hałas,
jest z~pewnością wielkim sukcesem i~zarazem potwierdzeniem, że nad tego typu
układami można pracować (z~różnymi rezultatami) korzystając z~amatorskich,
ogólnodostępnych i~tanich (lub darmowych) rozwiązań. \par Analizując rezultaty
przeprowadzonych badań od razu można stwierdzić, dlaczego wykorzystanie systemu
operacyjnego, który nie jest systemem czasu rzeczywistego, w przypadku
projektowania tego typu układów nie jest dobrą i~popularną decyzją. Opóźnienia
wprowadzane przez taki system i~trudności związane z~synchronizacją peryferiów
objawiają się mocno negatywnym wpływem na jakość działania finalnego układu. Jak
pokazano w sekcjach badającej wpływ wykorzystanej łatki PREEMPT\_RT, polepsza
ona działanie układu w porównaniu z~klasycznym jądrem. Wykorzystanie takich
rozszerzeń nie powoduje jednak, że system bazujący na jądrze Linux staje się
automatycznie systemem czasu rzeczywistego. Opóźnienia, pomimo że zmniejszone
i~bardziej deterministyczne, dalej występują i~ich wpływ dyskredytuje działanie
systemu dla sygnałów losowych i~takich, które szybko zmieniają swoją
częstotliwość. Adaptacja w przypadku takich sygnałów reaguje zbyt wolno i~hałasy
nie są tłumione poprawnie. Dodatkowo, system ALSA, jako że~nie jest
przystosowany do realizacji krytycznych czasowo operacji wejścia/wyjścia,
z~racji swojej elastyczności nie jest idealnym wyborem przy realizacji systemów
ANC. Jego głównym zadaniem jest zapewnienie warstwy abstrakcji wykorzystywanej
do komunikacji z~peryferiami audio, umożliwiającej konfigurowanie urządzeń
z~poziomu przestrzeni użytkownika, więc tworzenie pętli regulacji z~jego
wykorzystaniem nie jest z~pewnością wyborem optymalnym. \par Pomimo problemów
opisanych wyżej, zdaniem autora istnieje kilka sposobów, których wykorzystanie
może polepszyć działanie systemu. Sterowniki jądra odpowiadające za transfery
DMA oraz obsługę określonych urządzeń audio można zmodyfikować, zmieniając
rozmiar fragmentu oraz częstotliwość próbkowania na mniejsze. Docelowo, jeżeli
udałoby się uzyskać $n_f = 1$ i~$f_s = 1$ kHz, pętla przetwarzania mogłaby
operować na jednej próbce, jak to ma miejsce w typowych implementacjach na
procesory sygnałowe i~mikrokontrolery. Rozmiar fragmentu o~wartości jednej ramki
spowodowałby opóźnienie o~jedną ramkę, które jest krytyczne dla jakości
działania algorytmu. Z~częstotliwością próbkowania również można iść w drugą
stronę, zwiększając ją wraz z~krotnością decymacji w pętli przetwarzania. Takie
działanie, poza zmniejszeniem opóźnień, może również pozwolić na implementację
prostszej części analogowej, zmniejszając wymagania na projektowane filtry. \par
Poza modyfikacjami systemowymi, program można rozszerzyć o~wspomnianą w~sekcji
analizy tematu metodę niwelowania dodatkowej pętli sprzężenia, jeżeli wymagane
jest, aby układ działał na hałasach o~wyższym natężeniu lub/i~z~ułożeniem
sprzętu, powodującym wpływanie sygnału $y(n)$ na $x(n)$. W przypadku
zaprojektowanej elektorniki, patrząc z~perspektywy czasu autor uważa, że wybrane
przetworniki mogłyby być mniejszej dokładności (przykładowo 12 bitowe), a~filtry
o~mniejszym stopniu, redukując koszta i upraszczając tory analogowe. \par
Finalnie, podsumowując cały projekt można stwierdzić, że skonstruowany system
działa zgodnie z przyjętymi założeniami. Każdy element dobrze ze sobą współgra,
tworząc wynikowo układ zdolny do tłumienia hałasów wąskopasmowych w
akceptowalnym czasie. Zaimplementowane środowisko programistyczne może zostać
wykorzystane do dalszych prac, ponieważ jest łatwe w obsłudze, skalowalne i
proste w modyfikacji, a skonstruowana elektronika bardzo dobrze odnajduje się
jako karta pomiarowa dla sygnałów audio.
